% =====================================================================
\chapter{Teoremas del valor medio y aplicaciones}
\label{tema:5}
% =====================================================================

La derivada es una información local: describe $f$ en un entorno infinitesimal de cada punto. Los teoremas de este tema convierten esa información local en información global sobre un intervalo: si la derivada es positiva en todo un intervalo, la función crece en él; si está acotada, la función no puede variar más que una cantidad proporcional a la longitud del intervalo. El instrumento es el teorema del valor medio, del que se derivan la regla de L'Hôpital, los criterios de extremos y la teoría de la convexidad. Se siguen el capítulo~11 de Spivak~\cite{spivak}, el capítulo~4 de Apostol~\cite{apostol}, la §5.3 de los apuntes de \emph{Analysis~I}~\cite{cambA1}, las §4.1--4.5 y~4.7 de Stewart~\cite{stewart} y la unidad~2 del curso 18.01SC del MIT~\cite{mit1801}.

\section{Extremos locales y el teorema de Rolle}

\begin{definicion}[Extremos]\label{def:extremos}
Sea $f\colon D\to\R$ y $c\in D$. Se dice que $f$ tiene en $c$ un \emph{máximo absoluto} (o global) si $f(x)\le f(c)$ para todo $x\in D$, y un \emph{máximo local} (o relativo) si existe $\delta>0$ con $f(x)\le f(c)$ para todo $x\in D$ con $\abs{x-c}<\delta$. Análogamente se definen los mínimos. Un punto $c$ interior a $D$ con $f'(c)=0$ se llama \emph{punto crítico} o estacionario.
\end{definicion}

\begin{teorema}[Fermat: condición necesaria de extremo]\label{teo:fermat}
Si $f$ tiene un extremo local en un punto $c$ interior a su dominio y es derivable en $c$, entonces $f'(c)=0$.
\end{teorema}

\begin{proof}
Supóngase que es un máximo local. Para $0<h<\delta$, $f(c+h)-f(c)\le0$, luego $\dfrac{f(c+h)-f(c)}{h}\le0$ y, en el límite $h\to0^{+}$, $f'(c)\le0$. Para $-\delta<h<0$, el cociente es $\ge0$ (numerador $\le0$, denominador $<0$) y, en el límite $h\to0^{-}$, $f'(c)\ge0$. Luego $f'(c)=0$. El caso del mínimo se obtiene aplicando lo anterior a $-f$.
\end{proof}

El recíproco es falso: $x^{3}$ tiene $f'(0)=0$ y no tiene extremo en~0. Tampoco se puede omitir la hipótesis de derivabilidad ($\abs x$ tiene un mínimo en~0) ni la de punto interior ($x$ en $[0,1]$ tiene su máximo en~1, con $f'(1)=1$).

\begin{teorema}[Rolle]\label{teo:rolle}
Sea $f$ continua en $[a,b]$ y derivable en $(a,b)$, con $f(a)=f(b)$. Entonces existe $c\in(a,b)$ con $f'(c)=0$.
\end{teorema}

\begin{proof}
Por el \cref{teo:weierstrass}, $f$ alcanza en $[a,b]$ un máximo $M$ y un mínimo $m$. Si ambos se alcanzan en los extremos, entonces $M=m=f(a)$, $f$ es constante y $f'(c)=0$ para todo $c\in(a,b)$. Si no, uno de ellos se alcanza en un punto $c\in(a,b)$, que es interior, y el \cref{teo:fermat} da $f'(c)=0$.
\end{proof}

\section{El teorema del valor medio}

\begin{teorema}[Teorema del valor medio de Lagrange]\label{teo:tvm}
Sea $f$ continua en $[a,b]$ y derivable en $(a,b)$. Entonces existe $c\in(a,b)$ con
\[
f'(c)=\frac{f(b)-f(a)}{b-a},\qquad\text{es decir,}\qquad f(b)-f(a)=f'(c)(b-a).
\]
\end{teorema}

\begin{proof}
Sea $m=\frac{f(b)-f(a)}{b-a}$, la pendiente de la cuerda, y $g(x)=f(x)-m(x-a)$, que es $f$ menos la recta de pendiente $m$. Entonces $g$ es continua en $[a,b]$, derivable en $(a,b)$, y $g(a)=f(a)$, $g(b)=f(b)-(f(b)-f(a))=f(a)$. Por el \cref{teo:rolle}, existe $c\in(a,b)$ con $0=g'(c)=f'(c)-m$.
\end{proof}

Geométricamente: hay un punto en el que la tangente es paralela a la cuerda. Cinemáticamente: si un coche recorre 121 kilómetros en una hora y cincuenta minutos, su velocidad media es de 66~km/h, y en algún instante el velocímetro marca exactamente esa velocidad.

\begin{corolario}[Consecuencias del teorema del valor medio]\label{cor:tvm}
Sea $f$ continua en un intervalo $I$ y derivable en su interior.
\begin{enumerate}
  \item Si $f'=0$ en el interior de $I$, $f$ es constante en $I$.
  \item Si $f'=g'$ en el interior de $I$, entonces $f-g$ es constante.
  \item Si $f'>0$ (respectivamente $f'\ge0$) en el interior de $I$, $f$ es estrictamente creciente (respectivamente creciente) en $I$; análogamente con $<0$ y $\le0$.
  \item Si $\abs{f'(x)}\le M$ en el interior de $I$, entonces $\abs{f(x)-f(y)}\le M\abs{x-y}$ para todos $x,y\in I$; se dice que $f$ es \emph{lipschitziana} de constante~$M$.
\end{enumerate}
\end{corolario}

\begin{proof}
Sean $x<y$ en $I$. Por el \cref{teo:tvm} en $[x,y]$, $f(y)-f(x)=f'(c)(y-x)$ para algún $c\in(x,y)$. (i) $f'(c)=0$ da $f(y)=f(x)$. (ii) Es (i) aplicado a $f-g$. (iii) $f'(c)>0$ y $y-x>0$ dan $f(y)>f(x)$. (iv) $\abs{f(y)-f(x)}=\abs{f'(c)}\abs{y-x}\le M\abs{y-x}$.
\end{proof}

\begin{ejemplo}[Desigualdades con el teorema del valor medio]\label{ej:desigualdades}
(a) $\abs{\sin x-\sin y}\le\abs{x-y}$ para todos $x,y$, porque $\abs{\sin'}=\abs{\cos}\le1$.
(b) $e^{x}\ge1+x$ para todo $x\in\R$, con igualdad solo en $x=0$. Si $x>0$, el \cref{teo:tvm} en $[0,x]$ da $e^{x}-1=e^{c}x$ con $c>0$, luego $e^{x}-1>x$. Si $x<0$, en $[x,0]$: $1-e^{x}=e^{c}(-x)$ con $c<0$, luego $1-e^{x}<-x$.
(c) $\ln(1+x)\le x$ para $x>-1$: es (b) con $x$ sustituido por $\ln(1+x)$.
(d) $\sigma$ es lipschitziana de constante $\tfrac14$ y $\tanh$ de constante~1 (\cref{prop:derivadas-activaciones}).
\end{ejemplo}

\begin{teorema}[Teorema del valor medio de Cauchy]\label{teo:cauchy}
Sean $f$ y $g$ continuas en $[a,b]$ y derivables en $(a,b)$, con $g'(x)\ne0$ para todo $x\in(a,b)$. Entonces $g(b)\ne g(a)$ y existe $c\in(a,b)$ con
\[
\frac{f'(c)}{g'(c)}=\frac{f(b)-f(a)}{g(b)-g(a)}.
\]
\end{teorema}

\begin{proof}
Si fuera $g(b)=g(a)$, Rolle daría un punto con $g'=0$. Sea $k=\frac{f(b)-f(a)}{g(b)-g(a)}$ y $h(x)=f(x)-kg(x)$. Entonces $h(a)=h(b)$ (\cref{eje:5-cauchy}) y, por Rolle, existe $c\in(a,b)$ con $0=h'(c)=f'(c)-kg'(c)$.
\end{proof}

\section{La regla de L'Hôpital}

\begin{teorema}[Regla de L'Hôpital]\label{teo:lhopital}
Sean $f$ y $g$ derivables en un intervalo abierto que contiene a $a$, salvo quizá en $a$, con $g'(x)\ne0$ para $x\ne a$, y supóngase que $\lim_{x\to a}f(x)=\lim_{x\to a}g(x)=0$. Si existe $\lim_{x\to a}\dfrac{f'(x)}{g'(x)}=\ell$ (finito o infinito), entonces
\[
\lim_{x\to a}\frac{f(x)}{g(x)}=\ell.
\]
El resultado vale también para límites laterales, para $a=\pm\infty$, y cuando $\lim f=\lim g=\pm\infty$ (forma $\infty/\infty$).
\end{teorema}

\begin{proof}
Se prueba el caso $0/0$ con $a$ finito y $x\to a^{+}$; el caso $x\to a^{-}$ es simétrico. Se redefine (o define) $f(a)=g(a)=0$, con lo que $f$ y $g$ son continuas en $[a,x]$ para $x>a$ próximo. Por el \cref{teo:cauchy} en $[a,x]$, existe $c_x\in(a,x)$ con
\[
\frac{f(x)}{g(x)}=\frac{f(x)-f(a)}{g(x)-g(a)}=\frac{f'(c_x)}{g'(c_x)}.
\]
Cuando $x\to a^{+}$, también $c_x\to a^{+}$ (pues $a<c_x<x$), luego el último cociente tiende a $\ell$. La demostración de los demás casos puede verse en~\cite[cap.~11]{spivak}.
\end{proof}

\begin{ejemplo}[Aplicaciones de L'Hôpital]\label{ej:lhopital}
(a) $\lim_{x\to0}\dfrac{x-\sin x}{x^{3}}\overset{0/0}{=}\lim\dfrac{1-\cos x}{3x^{2}}=\dfrac16$, por la \cref{prop:senx}.
(b) $\lim_{x\to+\infty}\dfrac{\ln x}{x^{p}}\overset{\infty/\infty}{=}\lim\dfrac{1/x}{px^{p-1}}=\lim\dfrac{1}{px^{p}}=0$ para $p>0$.
(c) Forma $0\cdot\infty$: $\lim_{x\to0^{+}}x\ln x=\lim\dfrac{\ln x}{1/x}\overset{\infty/\infty}{=}\lim\dfrac{1/x}{-1/x^{2}}=\lim(-x)=0$.
(d) Forma $0^{0}$: $x^{x}=e^{x\ln x}\to e^{0}=1$ cuando $x\to0^{+}$, por (c).
(e) Forma $1^{\infty}$: $\lim_{x\to0}(1+\sin x)^{1/x}=\exp\Bigl(\lim\dfrac{\ln(1+\sin x)}{x}\Bigr)$, y por L'Hôpital el exponente tiende a $\lim\frac{\cos x}{1+\sin x}=1$; el límite es $e$.
\end{ejemplo}

\begin{proposicion}[Jerarquía de crecimientos]\label{prop:crecimientos}
Para todos $p>0$ y $b>1$: $\displaystyle\lim_{x\to+\infty}\frac{\ln x}{x^{p}}=0$ y $\displaystyle\lim_{x\to+\infty}\frac{x^{p}}{b^{x}}=0$.
\end{proposicion}

\begin{proof}
La primera es el \cref{ej:lhopital}(b). Para la segunda, sea $n\ge p$ entero. Aplicando L'Hôpital $n$ veces, $\lim\frac{x^{n}}{b^{x}}=\lim\frac{n!}{b^{x}(\ln b)^{n}}=0$, y $0<\frac{x^{p}}{b^{x}}\le\frac{x^{n}}{b^{x}}$ para $x\ge1$.
\end{proof}

\section{Monotonía, extremos y optimización}

\begin{proposicion}[Criterios de extremo local]\label{prop:criterios}
Sea $f$ derivable en un entorno de un punto crítico $c$.
\begin{enumerate}
  \item (\emph{Criterio de la primera derivada}) Si $f'\ge0$ a la izquierda de $c$ y $f'\le0$ a la derecha (en un entorno), $f$ tiene un máximo local en $c$; con los signos opuestos, un mínimo local.
  \item (\emph{Criterio de la segunda derivada}) Si $f''(c)$ existe y $f''(c)<0$, $f$ tiene un máximo local estricto en $c$; si $f''(c)>0$, un mínimo local estricto. Si $f''(c)=0$, el criterio no decide.
\end{enumerate}
\end{proposicion}

\begin{proof}
(i) Por el \cref{cor:tvm}(iii), $f$ es creciente en $(c-\delta,c]$ y decreciente en $[c,c+\delta)$, luego $f(x)\le f(c)$ en $(c-\delta,c+\delta)$. (ii) Si $f''(c)<0$, entonces $\lim_{h\to0}\frac{f'(c+h)-f'(c)}{h}=\lim\frac{f'(c+h)}{h}<0$, y para $h$ pequeño $f'(c+h)/h<0$: $f'>0$ a la izquierda de $c$ y $f'<0$ a la derecha. Por la demostración de (i), con monotonía estricta, $c$ es un máximo local estricto. Los ejemplos $x^{4}$, $-x^{4}$ y $x^{3}$ en $c=0$ muestran que con $f''(c)=0$ puede haber mínimo, máximo o ninguno de los dos.
\end{proof}

\begin{proposicion}[Extremos absolutos en un intervalo cerrado]\label{prop:intervalo-cerrado}
Sea $f$ continua en $[a,b]$. Su máximo y su mínimo absolutos (que existen por el \cref{teo:weierstrass}) se alcanzan en los extremos $a$, $b$, en puntos críticos de $(a,b)$ o en puntos de $(a,b)$ donde $f$ no es derivable. Basta, pues, evaluar $f$ en esos candidatos y comparar.
\end{proposicion}

\begin{proof}
Si el máximo se alcanza en un punto $c\in(a,b)$ en el que $f$ es derivable, el \cref{teo:fermat} da $f'(c)=0$.
\end{proof}

\begin{ejemplo}[Un problema de optimización]\label{ej:optimizacion}
Se quiere construir una caja abierta recortando cuadrados de lado $x$ en las esquinas de una lámina cuadrada de lado~12 y doblando los bordes. El volumen es $V(x)=x(12-2x)^{2}$ para $x\in[0,6]$. Entonces $V'(x)=(12-2x)^{2}-4x(12-2x)=(12-2x)(12-6x)$, que se anula en $x=2$ y $x=6$. Los candidatos dan $V(0)=V(6)=0$ y $V(2)=2\cdot64=128$: el volumen máximo es $128$, con $x=2$.
\end{ejemplo}

\section{Convexidad}

\begin{definicion}[Función convexa]\label{def:convexa}
Una función $f\colon I\to\R$ en un intervalo $I$ es \emph{convexa} si para todos $x,y\in I$ y todo $t\in[0,1]$
\[
f\bigl(tx+(1-t)y\bigr)\le tf(x)+(1-t)f(y),
\]
es decir, si el segmento que une dos puntos de la gráfica queda por encima de ella. Es \emph{estrictamente convexa} si la desigualdad es estricta para $x\ne y$ y $t\in(0,1)$, y \emph{cóncava} si $-f$ es convexa. Un \emph{punto de inflexión} es un punto donde $f$ es continua y pasa de convexa a cóncava o al revés.
\end{definicion}

\begin{teorema}[Caracterizaciones de la convexidad]\label{teo:convexidad}
Sea $f$ derivable en un intervalo abierto $I$.
\begin{enumerate}
  \item Si $f'$ es creciente en $I$, entonces $f$ está por encima de sus tangentes: $f(y)\ge f(x)+f'(x)(y-x)$ para todos $x,y\in I$.
  \item Si $f$ está por encima de sus tangentes, $f$ es convexa.
  \item En particular, si $f''\ge0$ en $I$, $f$ es convexa; si $f''>0$, es estrictamente convexa.
\end{enumerate}
\end{teorema}

\begin{proof}
(i) Sea $y>x$. Por el \cref{teo:tvm}, $f(y)-f(x)=f'(c)(y-x)$ con $c\in(x,y)$, y $f'(c)\ge f'(x)$ por el crecimiento de $f'$; como $y-x>0$, $f(y)-f(x)\ge f'(x)(y-x)$. Si $y<x$, se aplica el TVM en $[y,x]$: $f(x)-f(y)=f'(c)(x-y)\le f'(x)(x-y)$, que es la misma desigualdad.
(ii) Sean $x,y\in I$, $t\in[0,1]$ y $z=tx+(1-t)y$. Las desigualdades de la tangente en $z$ dan $f(x)\ge f(z)+f'(z)(x-z)$ y $f(y)\ge f(z)+f'(z)(y-z)$. Multiplicando por $t$ y por $1-t$ y sumando, $tf(x)+(1-t)f(y)\ge f(z)+f'(z)\bigl(tx+(1-t)y-z\bigr)=f(z)$.
(iii) Si $f''\ge0$, $f'$ es creciente por el \cref{cor:tvm}(iii), y se aplican (i) y (ii). La versión estricta se obtiene igual con desigualdades estrictas.
\end{proof}

\begin{ejemplo}[Convexidad de funciones elementales y de activación]\label{ej:convexas}
$x^{2}$, $e^{x}$ y $-\ln x$ son estrictamente convexas ($f''=2$, $e^{x}$ y $1/x^{2}$). La softplus es estrictamente convexa: $\softplus''=\sigma'>0$. La sigmoide es convexa en $(-\infty,0]$ y cóncava en $[0,\infty)$, porque $\sigma''=\sigma(1-\sigma)(1-2\sigma)$ tiene el signo de $1-2\sigma$; su único punto de inflexión es $x=0$. La $\ReLU$ es convexa (es el máximo de dos funciones afines, \cref{eje:5-max-convexas}), aunque no sea derivable.
\end{ejemplo}

\begin{observacion}[Errores frecuentes]\label{obs:5-errores}
(a) Aplicar L'Hôpital sin comprobar que la forma es $0/0$ o $\infty/\infty$: $\lim_{x\to0}\frac{\cos x}{x+1}=1$, mientras que el cociente de derivadas tiende a~$0$.
(b) Derivar el cociente con la regla del cociente en vez de derivar numerador y denominador por separado.
(c) Concluir que si $\lim f'/g'$ no existe, tampoco existe $\lim f/g$: $\frac{x+\sin x}{x}\to1$ cuando $x\to\infty$, aunque $\frac{1+\cos x}{1}$ no tenga límite.
(d) Identificar «punto crítico» con «extremo», o buscar extremos absolutos solo entre los puntos críticos, olvidando los extremos del intervalo y los puntos sin derivada.
(e) Pensar que $f''(c)=0$ implica punto de inflexión: $x^{4}$ tiene $f''(0)=0$ y es convexa en todo $\R$.
\end{observacion}

\begin{observacion}[Conexión con el curso]
El \cref{cor:tvm}(iv) es la idea del módulo~7: si el gradiente está acotado (o es lipschitziano), un paso pequeño del descenso de gradiente cambia poco la pérdida, y el \emph{lema de descenso} cuantifica cuánto disminuye. La convexidad garantiza que todo mínimo local es global (módulos~4 y~7), y la desigualdad de la tangente $e^{x}\ge1+x$, o su forma $\ln t\le t-1$, es el paso clave de la desigualdad de Gibbs del módulo~6. La regla de L'Hôpital justifica el convenio $0\ln0=0$ en la entropía ($x\ln x\to0$, \cref{ej:lhopital}(c)), y la \cref{prop:crecimientos} es el lenguaje con el que se comparan los costes cuadrático y lineal de la atención o las leyes de potencia del módulo~16.
\end{observacion}

\begin{lecturas}
\item Spivak~\cite{spivak}, cap.~11 «Significance of the Derivative».
\item Apostol~\cite{apostol}, cap.~4 «Differential Calculus».
\item Gowers y Chua~\cite{cambA1}, §5.3 «Differentiation theorems».
\item Stewart~\cite{stewart}, §4.1--4.5 y §4.7; OpenStax~\cite{osc1}, §4.3--4.5, §4.7 y §4.8.
\item MIT 18.01SC~\cite{mit1801}, unidad~2: sesiones 27--30 (gráficas y optimización) y 34--35 (teorema del valor medio).
\end{lecturas}

\section*{Ejercicios del tema 5}
\addcontentsline{toc}{section}{Ejercicios del tema 5}

\begin{ejercicio}{1}{eje:5-hipotesis}
Para cada función, diga si se cumplen las hipótesis del teorema del valor medio en el intervalo dado y, si no, si aun así existe algún $c$ con $f'(c)=\frac{f(b)-f(a)}{b-a}$: (a) $\abs{x-\tfrac12}$ en $[-1,1]$; (b) $1/x^{2}$ en $[-1,1]$; (c) $e^{x}$ en $[0,1]$; (d) $\sqrt[3]{x}$ en $[-1,1]$.
\fuente{adaptado de OpenStax~\cite{osc1}, §4.4, ejercicios de la sección.}
\end{ejercicio}

\begin{ejercicio}{2}{eje:5-valores-c}
Halle todos los puntos $c$ que proporciona el teorema del valor medio para: (a) $x^{2}$ en $[0,1]$; (b) $\ln x$ en $[1,2]$; (c) $x^{3}-x$ en $[-2,2]$; (d) $x^{3}$ en $[0,2]$.
\fuente{adaptado de MIT 18.01SC, Jerison y Mattuck~\cite{mit1801ex}, ejercicio 2G-1, y OpenStax~\cite{osc1}, §4.4.}
\end{ejercicio}

\begin{ejercicio}{2}{eje:5-lhopital}
Calcule: (a) $\lim_{x\to0}\dfrac{e^{x}-x-1}{x^{2}}$; (b) $\lim_{x\to0}\dfrac{\sin x-\tan x}{x^{3}}$; (c) $\lim_{x\to1}\dfrac{x-1}{\ln x}$; (d) $\lim_{x\to\infty}x^{2}e^{-x}$; (e) $\lim_{x\to0^{+}}x^{2x}$; (f) $\lim_{x\to\infty}\Bigl(1-\dfrac1x\Bigr)^{x}$; (g) $\lim_{x\to0}\Bigl(\dfrac{1}{\sin x}-\dfrac1x\Bigr)$; (h) $\lim_{x\to0}\dfrac{3^{x}-2^{x}}{x}$.
\fuente{adaptado de OpenStax~\cite{osc1}, §4.8, ejercicios de la sección.}
\end{ejercicio}

\begin{ejercicio}{2}{eje:5-grafica}
Estudie dominio, asíntotas, intervalos de crecimiento, extremos locales, convexidad y puntos de inflexión, y esboce la gráfica de: (a) $f(x)=\dfrac{x^{2}}{x-1}$; (b) $g(x)=x-3x^{1/3}$; (c) $h(x)=e^{-x^{2}}$.
\fuente{adaptado de MIT 18.01SC, Jerison y Mattuck~\cite{mit1801ex}, ejercicios 2B-1(d) y 2B-1(h), y examen~2, problema~3.}
\end{ejercicio}

\begin{ejercicio}{2}{eje:5-optimizacion}
(a) Maximice $xy$ sujeto a $x+2y=a$, con $a>0$. (b) Halle el máximo de $x^{2}\ln(1/x)$ en $(0,\infty)$. (c) Estudie si $(x+1)e^{-x}$ alcanza un máximo y un mínimo en $[0,\infty)$. (d) Halle el rectángulo de área máxima inscrito en una semicircunferencia de radio $a$, con un lado sobre el diámetro.
\fuente{adaptado de MIT 18.01SC, Jerison y Mattuck~\cite{mit1801ex}, ejercicios 2C-3, 2C-7, 2C-14(a) y 2C-15.}
\end{ejercicio}

\begin{ejercicio}{2}{eje:5-cubica}
Sea $y=x^{3}+ax^{2}+bx+c$. Halle una condición necesaria y suficiente sobre $a$, $b$, $c$ para que $y$ tenga un máximo local y un mínimo local. Compruébela con $x^{3}-3x+1$ y con $x^{3}+x$.
\fuente{adaptado de MIT 18.01SC, Jerison y Mattuck~\cite{mit1801ex}, ejercicio 2B-3.}
\end{ejercicio}

\begin{ejercicio}{2}{eje:5-una-raiz}
Demuestre que la ecuación $x^{3}+3x^{2}+16=0$ tiene exactamente una raíz real, y hállela. Demuestre también que $e^{x}=2-x$ tiene exactamente una solución real.
\fuente{adaptado de OpenStax~\cite{osc1}, §4.4, ejercicio de la sección sobre $x^{3}+3x^{2}+16$.}
\end{ejercicio}

\begin{ejercicio}{3}{eje:5-desigualdades}
Demuestre con el teorema del valor medio: (a) $\sin x<x$ para $x>0$; (b) $\sqrt{1+x}<1+\dfrac x2$ para $x>0$; (c) $\dfrac{x}{1+x}<\ln(1+x)<x$ para $x>0$; (d) $\abs{\arctan x-\arctan y}\le\abs{x-y}$.
\fuente{adaptado de MIT 18.01SC, Jerison y Mattuck~\cite{mit1801ex}, ejercicio 2G-2, y examen~2, problema~6.}
\end{ejercicio}

\begin{ejercicio}{3}{eje:5-comparacion}
(a) Demuestre que si $f$ y $g$ son derivables en $[a,\infty)$, $f(a)\ge g(a)$ y $f'(x)>g'(x)$ para $x>a$, entonces $f(x)>g(x)$ para $x>a$. (b) Aplíquelo sucesivamente para demostrar que, para $x>0$, $e^{x}>1+x$ y $e^{x}>1+x+\dfrac{x^{2}}{2}$. (c) Demuestre por inducción que $e^{x}>\sum_{k=0}^{n}x^{k}/k!$ para $x>0$ y todo $n$.
\fuente{adaptado de MIT 18.01SC, Jerison y Mattuck~\cite{mit1801ex}, ejercicio 2G-7.}
\end{ejercicio}

\begin{ejercicio}{3}{eje:5-raices-polinomio}
(a) Demuestre que si un polinomio de grado $n$ tiene $n$ raíces reales distintas, su derivada tiene $n-1$ raíces reales distintas. (b) Sea $f$ dos veces derivable con ceros en $a<b<c$; demuestre que existe $p\in(a,c)$ con $f''(p)=0$. (c) Demuestre que $x^{5}+x-1$ tiene exactamente una raíz real.
\fuente{adaptado de MIT 18.01SC, Jerison y Mattuck~\cite{mit1801ex}, ejercicios 2G-4 y 2G-5.}
\end{ejercicio}

\begin{ejercicio}{3}{eje:5-tvm-demo}
(a) Enuncie el teorema de Rolle y demuestre a partir de él el teorema del valor medio. (b) Sea $f$ dos veces derivable en $\R$. Demuestre que
\[
f''(a)=\lim_{h\to0}\frac{f(a+h)+f(a-h)-2f(a)}{h^{2}}.
\]
(Indicación: aplique L'Hôpital una vez y use la definición de $f''(a)$.)
\fuente{adaptado de Cambridge, \emph{Tripos} IA 2013, Paper~1~\cite{tripos2013}, pregunta~11E(i), (ii) y (iv).}
\end{ejercicio}

\begin{ejercicio}{3}{eje:5-cauchy}
(a) Complete la demostración del \cref{teo:cauchy} comprobando que $h(a)=h(b)$. (b) Interprete el teorema para la curva paramétrica $t\mapsto(g(t),f(t))$. (c) Deduzca la regla de L'Hôpital del caso $0/0$ con $x\to a^{+}$ y explique dónde se usa la hipótesis $g'\ne0$.
\fuente{adaptado de Cambridge, \emph{Tripos} IA 2013, Paper~1~\cite{tripos2013}, pregunta~11E(iii).}
\end{ejercicio}

\begin{ejercicio}{3}{eje:5-darboux}
(Teorema de Darboux.) Sea $g$ derivable en un intervalo abierto que contiene a $[a,b]$, con $g'(a)<k<g'(b)$. Demuestre que existe $c\in(a,b)$ con $g'(c)=k$. (Indicación: la función $\varphi(x)=g(x)-kx$ es continua en $[a,b]$; demuestre que su mínimo no se alcanza en $a$ ni en $b$.) Deduzca que una derivada no puede tener discontinuidades de salto y que la función escalón no es la derivada de ninguna función.
\fuente{adaptado de Cambridge, \emph{Tripos} IA 2019, Paper~1~\cite{tripos2019}, pregunta~10D.}
\end{ejercicio}

\begin{ejercicio}{3}{eje:5-creciente-derivada}
(a) Demuestre que si $f$ es creciente en un intervalo abierto y derivable en un punto $a$ de él, entonces $f'(a)\ge0$. (b) Explique por qué el argumento no da $f'(a)>0$ aunque $f$ sea estrictamente creciente, y dé un contraejemplo. (c) Demuestre que $\sigma$ y $\tanh$ son estrictamente crecientes y que $\ReLU$ es creciente pero no estrictamente.
\fuente{adaptado de MIT 18.01SC, Jerison y Mattuck~\cite{mit1801ex}, ejercicio 2B-7.}
\end{ejercicio}

\begin{ejercicio}{3}{eje:5-max-convexas}
(a) Demuestre con la \cref{def:convexa} que el máximo de dos funciones convexas es convexo, y deduzca que $\ReLU$ y $\abs{x}$ son convexas. (b) Demuestre que la suma de funciones convexas es convexa y que, si $f$ es convexa y $g$ es convexa y creciente, $g\circ f$ es convexa. (c) Deduzca que $x\mapsto\softplus(x^{2}-1)$ es convexa.
\end{ejercicio}

\begin{ejercicio}{3}{eje:5-jensen}
(a) Demuestre por inducción la desigualdad de Jensen finita: si $f$ es convexa, $t_1,\dots,t_n\ge0$ con $\sum t_i=1$, entonces $f\bigl(\sum t_ix_i\bigr)\le\sum t_if(x_i)$. (b) Aplíquela a $-\ln$ para obtener la desigualdad entre las medias aritmética y geométrica: $\sqrt[n]{x_1\cdots x_n}\le\frac{x_1+\dots+x_n}{n}$ para $x_i>0$.
\end{ejercicio}
